\section*{Bab \ref{ch:b3:hydrogeology} --- Hidrogeologi}

\begin{solution}{exo:b3:hydrogeology:1}
Pada layar pertama: $z = 120 - 40 = \qty{80}{m}$, $h = 120 - 12.5 =
\qty{107.5}{m}$, tinggi tekanan $107.5 - 80 = \qty{27.5}{m}$. Pada layar
kedua: $z = \qty{40}{m}$, $h = \qty{105.0}{m}$, tinggi tekanan
\qty{65.0}{m}. \omterm{def:b3:hydrogeology:head}{Tinggi hidraulik} lebih besar di layar dangkal, sehingga air
mengalir ke bawah, di bawah gradien $(107.5 - 105.0)/40 = \num{0.0625}$:
piezometer berada di daerah imbuhan. Tekanan saja
akan menyarankan sebaliknya.
\end{solution}

\begin{solution}{exo:b3:hydrogeology:2}
$q = K\Delta h/L = 12\times0.30/0.50 = \qty{7.2}{m/d}$. Luasnya
$\pi\times0.05^2 = \qty{7.85e-3}{m^2}$, sehingga $Q = \qty{0.0565}{m^3/d}$, atau
$56.5/24 = \qty{2.36}{L/h}$. \omterm{def:b3:hydrogeology:seepage}{Kecepatan rembesannya} $7.2/0.35 =
\qty{20.6}{m/d}$.
\end{solution}

\begin{solution}{exo:b3:hydrogeology:3}
$K = 2.5/86400 = \qty{2.894e-5}{m/s}$, dan $k = K\mu/(\rho g) =
\num{2.894e-5}\times\num{1.306e-3}/(999.7\times9.81) =
\qty{3.85e-12}{m^2}$. Pada \qty{20}{\celsius}, $K$ dikalikan
$(998.2/\num{1.002e-3})/(999.7/\num{1.306e-3}) = \num{1.30}$: $K =
\qty{3.25}{m/d}$. Perubahannya hampir seluruhnya berasal dari viskositas.
\end{solution}

\begin{solution}{exo:b3:hydrogeology:4}
$T = 2.5\times180 = \qty{450}{m^2/d}$. Ekspansi air memberikan
$\num{9e-7}\times180 = \num{1.62e-4}$, sekitar \qty{80}{\percent} dari $S =
\num{2e-4}$; kompresi batu gamping menyumbang sisanya.
\end{solution}

\begin{solution}{exo:b3:hydrogeology:5}
Menurut \cref{prop:b3:hydrogeology:thiem}, $T = Q\ln(r_2/r_1)/(2\pi(s_1 -
s_2)) = 1500\times\ln10/(2\pi\times1.53) = \qty{359}{m^2/d}$, dan $K =
359/40 = \qty{9.0}{m/d}$.
\end{solution}

\begin{solution}{exo:b3:hydrogeology:6}
$Q/4\pi T = 1200/(4\pi\times300) = \qty{0.318}{m}$. Setelah 1 hari, $u =
300^2\times\num{5e-4}/(4\times300\times1) = \num{0.0375}$ dan $W = -0.5772
+ 3.2834 + 0.0375 - 0.0004 = \num{2.743}$: $s = \qty{0.873}{m}$. Setelah 10
hari, $u = \num{0.00375}$, $W = -0.5772 + 5.5860 + 0.0038 = \num{5.013}$:
$s = \qty{1.596}{m}$. Sepuluh kali lebih lama menghasilkan
penurunan muka air kurang dari dua kali lipat.
\end{solution}

\begin{solution}{exo:b3:hydrogeology:7}
$u \le 0.05$ ketika $t \ge r^2S/(4T\times0.05)$. Pada \qty{60}{m}: $3600\times
\num{2e-4}/90 = \qty{0.008}{d}$, yaitu \qty{11.5}{min}. Pada \qty{500}{m}:
\qty{0.556}{d}, \qty{800}{min}, lebih dari 13 jam. Waktunya tumbuh sebanding dengan $r^2$:
dalam uji satu atau dua hari, sumur pengamatan yang dipakai untuk garis lurus
harus berada dalam beberapa ratus meter dari sumur pompa.
\end{solution}

\begin{solution}{exo:b3:hydrogeology:8}
Di tengah, setiap sumur berjarak \qty{200}{m}: $u = 200^2\times\num{2e-4}/(4
\times450\times10) = \num{4.44e-4}$, $W = \num{7.142}$, $s = 0.382\times
7.142 = \qty{2.728}{m}$ dari masing-masing, total \qty{5.456}{m}. Pada
\qty{400}{m}, $u = \num{1.78e-3}$, $W = \num{5.757}$, $s =
\qty{2.199}{m}$. Setiap sumur menurunkan muka air di sumur lainnya sebesar
\qty{2.2}{m}: pompanya harus mengangkat air \qty{2.2}{m} lebih tinggi, yang
membutuhkan energi, dan penurunan yang tersisa sebelum intake pompa kering
menjadi lebih kecil.
\end{solution}

\begin{solution}{exo:b3:hydrogeology:9}
Empat garis aliran di antara dasar bendungan dan lantai membatasi $N_f = 5$
saluran; sembilan garis ekuipotensial membagi \qty{10}{m} menjadi $N_d = 10$
tetesan. $q = KHN_f/N_d = 5\times10\times5/10 = \qty{25}{m^3/d}$ per meter
bendungan, dan \qty{7500}{m^3/d} di bawah \qty{300}{m}. Faktor bentuk yang
dihitung, 0.512 alih-alih 0.5, akan memberikan \qty{25.6}{m^3/d} per meter.
\end{solution}

\begin{solution}{exo:b3:hydrogeology:10}
$k = na^2/8 = 0.30\times(\num{5e-5})^2/8 = \qty{9.4e-11}{m^2}$, dan $K =
k\rho g/\mu = \num{9.375e-11}\times998.2\times9.81/\num{1.002e-3} =
\qty{9.2e-4}{m/s}$, \qty{79}{m/d}, berada dalam rentang pasir bersih, 0.3 sampai
\qty{150}{m/d}. Model ini mengabaikan kelokan jalur, yang lebih panjang
daripada tabung lurus (ini menurunkan $K$), tenggorokan antarpori,
yang lebih sempit daripada pori (lebih rendah), pori yang berujung buntu dan tidak
mengalirkan aliran (lebih rendah), serta sebaran ukuran pori: karena tabung membawa
aliran sebanding dengan $a^4$, jalur terlebar membawa sebagian besar air dan
sebaran itu menaikkan $K$ di atas nilai tabung berukuran rata-rata.
\end{solution}

\begin{solution}{exo:b3:hydrogeology:11}
(a)~Logaritma menjadi nol ketika $2.25\,Tt = r^2S$, yaitu pada $r_0 =
\sqrt{2.25\,Tt/S}$. (b)~Terkekang dan tak terkekang:
\[
r_0 = \sqrt{\frac{2.25\times450\times30}{\num{2e-4}}} = \qty{12.3}{km},
\qquad r_0 = \sqrt{\frac{2.25\times450\times30}{0.2}} = \qty{390}{m},
\]
$\sqrt{1000} = 31.6$ kali lebih kecil. (c)~Tekanan
menyebar dengan difusivitas $T/S$, sangat besar ketika $S$ kecil: selama
uji, kerucut mencapai sesar, singkapan akuifer dan sungai,
dan penurunan muka air keluar dari garis lurus, yang memberi tahu
ahli hidrogeologi bahwa sebuah batas telah tercapai.
\end{solution}

\begin{solution}{exo:b3:hydrogeology:12}
(a)~Dilihat dari kerangka yang bergerak bersama air, front tersebar dengan
simpangan baku $\sigma_x = \sqrt{2Dt} = \sqrt{2\alpha_Lvt}$. Front melewati
$x$ sekitar $t = x/v$, ketika $\sigma_x = \sqrt{2\alpha_Lx}$; pada kecepatan
$v$, sebaran $\sigma_x$ dalam jarak setara dengan $\sigma_t = \sigma_x/v$
dalam waktu. (b)~Waktu \qty{16}{\percent} dan \qty{84}{\percent} berada satu
simpangan baku di kedua sisi rata-rata: $\sigma_t = (247.2 -
152.6)/2 = \num{47.3}\text{ hari}$. Maka $\alpha_L = v^2\sigma_t^2/(2x) =
0.25\times47.3^2/200 = \qty{2.8}{m}$. (Kurva dihitung dengan
\qty{3}{m}: aproksimasi normal sedikit meremehkannya.) (c)~$R =
1 + 1.9\times0.32/0.30 = \num{3.03}$: zat terlarut bergerak dengan sepertiga
kecepatan air.
\end{solution}

\begin{solution}{pb:b3:hydrogeology:1}
\textbf{1.} Serpih $363.2 - 282.5 = \qty{80.7}{m}$; batu gamping $282.5 -
102.5 = \qty{180}{m}$.

\textbf{2.} $h = 363.2 - 45.0 = \qty{318.2}{m}$; tinggi tekanan di puncak
$318.2 - 282.5 = \qty{35.7}{m}$. Air naik \qty{35.7}{m} di atas
puncak batu gamping, ditahan oleh serpih: akuifernya tertekan.

\textbf{3.} Dari $h_2 - h_1 = 550a - 100b = -1.62$ dan $h_3 - h_1 = 250a
- 600b = -1.44$: $a = \num{-2.715e-3}$ dan $b = \num{1.269e-3}$ per
meter.

\textbf{4.} $|\nabla h| = \sqrt{a^2 + b^2} = \num{0.0030}$. Air mengalir
searah $-\nabla h = (\num{2.715e-3}, \num{-1.269e-3})$, ke timur dan
selatan: azimut $\ang{180} - \arctan(2.715/1.269) = \ang{115}$.

\textbf{5.} Titik-titik sebelum sekitar \qty{10}{min} berada di atas garis:
di sana $u > 0.05$ dan kurva Theis berada di atas aproksimasi logaritmiknya
($W > -\gamma - \ln u$, seperti ditunjukkan oleh gambar fungsi sumur).

\textbf{6.} Garis terbaik melalui titik-titik dari \qty{15}{min} ke atas
melewati \qty{1.75}{m} pada \qty{100}{min} dan \qty{2.63}{m} pada \qty{1000}{min}:
$\Delta s = \qty{0.87}{m}$ per siklus logaritma.

\textbf{7.} $T = 2.3\times2160/(4\pi\times0.872) = \qty{453}{m^2/d}$.

\textbf{8.} Garis turun \qty{0.87}{m} per siklus dari \qty{1.75}{m} pada
\qty{100}{min}: nol tercapai $1.75/0.872 = 2.01$ siklus lebih awal,
$t_0 = \qty{0.98}{min} = \qty{6.83e-4}{d}$.

\textbf{9.} $S = 2.25\,Tt_0/r^2 = 2.25\times453\times\num{6.83e-4}/3600 =
\num{1.9e-4}$, berada dalam kisaran \omterm{def:b3:hydrogeology:aquifer-types}{akuifer tertekan}.

\textbf{10.} Pada \qty{15}{min}, $u = 3600\times\num{1.9e-4}/(4\times453
\times0.0104) = \num{0.036}$, dan lebih kecil sesudahnya.

\textbf{11.} $K = 453/180 = \qty{2.5}{m/d}$, batu gamping yang retak.

\textbf{12.} $u = 3600\times\num{2e-4}/(4\times450\times2) =
\num{2.0e-4}$, $W = -0.5772 + 8.5172 + 0.0002 = \num{7.940}$, $s =
0.382\times7.940 = \qty{3.03}{m}$, sesuai catatan.

\textbf{13.} $u = 10^6\times\num{2e-4}/(4\times450\times30) =
\num{0.0037}$, $W = \num{5.02}$, $s = \qty{1.92}{m}$.

\textbf{14.} $r_0 = \sqrt{2.25\times450\times30/(\num{2e-4})} =
\qty{12.3}{km}$, lebih lebar daripada lebar wilayah itu. Kerucut telah
mencapai tepi batu gamping, singkapannya, dan sesar, yang mungkin
menempatkannya berhadapan dengan batuan yang kurang permeabel: akuifer tak terhingga
dari Theis tidak lagi menggambarkannya, dan hanya catatan penurunan muka air
yang akan menunjukkan apakah suatu batas memasok sumur atau menghentikan aliran.

\textbf{15.} $2160\times30 = \qty{64800}{m^3}$. Air itu berasal dari
penyimpanan: air yang terlepas akibat pemuaian air dan pemampatan batu
gamping, di seluruh area kerucut. Mengintegralkan $S\,s$ atas kerucut
menghasilkan kembali $Qt$.

\textbf{16.} $u = 10^6\times0.18/(4\times450\times30) = \num{3.33}$, $W
= \num{0.0086}$, $s = \qty{3}{mm}$: peternakan itu tidak akan merasakannya.
Pemompaan akan menarik air dari pori-pori dekat sumur, bukan dari tekanan
seluruh akuifer.

\textbf{17.} $q = 2.5\times0.0030 = \qty{0.0075}{m/d}$; $v = 0.0075/0.18
= \qty{0.042}{m/d}$, yaitu \qty{15}{m} per tahun.

\textbf{18.} $200/0.0417 = \num{4800}\text{ hari}$, \num{13.1}~tahun; untuk
polutan, \num{32.9}~tahun.

\textbf{19.} Melalui silinder berjari-jari $r$, debit $Q$ melintasi
luas $2\pi rb$, yang fraksi porinya adalah $n_e$: $\mathrm dr/\mathrm
dt = -Q/(2\pi rbn_e)$, dan mengintegralkan $2\pi bn_er\,\mathrm dr = -Q\,
\mathrm dt$ dari $r$ ke 0 menghasilkan $t = \pi n_ebr^2/Q$. Di sini $\pi\times0.18
\times180\times200^2/2160 = \num{1885}\text{ hari}$, \num{5.2}\text{ tahun} untuk
air dan $2.5\times5.2 = \num{12.9}\text{ tahun}$ untuk polutan: sumur
membawanya lebih dari dua kali lebih cepat.

\textbf{20.} $\alpha_L = \qty{20}{m}$ dan $\sqrt{2\times20\times200} =
\qty{89}{m}$. Konsentrasi mencapai \qty{16}{\percent} dari nilai
akhirnya \qty{89}{m} di depan front rata-rata: tingkat itu mencapai
sumur ketika front rata-rata masih berjarak \qty{89}{m}, pada $(111/200)^2\times12.9 = \num{4.0}\text{ tahun}$, sepertiga
dari waktu rata-rata. Jejak pertama penting bagi air minum, bukan
rata-ratanya.

\textbf{21.} Pendukung: sumur menangkap polutan, sehingga polutan tidak
dapat menyebar lebih jauh ke arah hilir, dan air dapat diolah di tempat
ia dipompa. Penentang: pemompaan menarik polutan ke sumur pasokan lebih
cepat daripada jika ia bergerak ke mana pun, dan air harus dipantau serta
diolah selama puluhan tahun.

\textbf{22.} \omterm{def:b3:hydrogeology:storage}{Transmisivitas} batu gamping Varden adalah $T \approx
\textbf{\qty[text-series-to-math=true]{450}{m^2/d}}$ (\qty{453}{m^2/d}
dari uji).
\end{solution}
